\documentclass[11pt]{article}
\usepackage{fontspec}
\setmainfont{TeX Gyre Pagella}
\usepackage{microtype}
\usepackage[margin=1.3in]{geometry}
\usepackage{amsmath,amssymb}
\usepackage[hidelinks]{hyperref}
\usepackage{enumitem}
\setlist{noitemsep}

\title{Epistemic Immutability: Provenance, Counterexamples, and Selective Continuation}
\author{Flyxion \\ \small Independent Researcher}
\date{August 2026}

\begin{document}
\maketitle

\begin{abstract}
An artifact's metadata records a claim about its history; it does not, by itself, establish that history. A file's modification date, a Git commit's timestamp, and a model's prediction about an unobserved state are all declarations made from within a limited evidential horizon, and all three can be preserved with perfect integrity while being false, or true only in a narrower sense than they appear to assert. This essay develops a single formal distinction, between declared history and witnessed history, and shows that a small number of mundane technical cases, an audio recorder with an uninitialized clock, a Git repository with rewritten commit dates, and a vision model asked to predict the unseen side of an object, are instances of the same underlying structure. The central proposal is that later evidence should constrain interpretation without rewriting the original claim: preserve the commitment, append the correction, record the discrepancy, and correct the continuation rather than the record. This principle is given a measure-theoretic formalization in terms of admissible interpretation spaces, and its constructive counterpart, the media quine, a provenance-bearing family of transformations of a single artifact, is sketched as an architecture that treats disagreement between representations as evidence rather than noise.
\end{abstract}

\section{Two Kinds of Claim About the Past}

A timestamp is not a report of an event; it is a property of a record. This distinction sounds pedantic until it is pressed on a concrete case. A digital audio recorder that requires a date to be set before use, and whose owner advances through that field quickly without attending to it, will faithfully attach a syntactically valid but historically arbitrary date to every subsequent recording. Nothing in this sequence involves malfunction. The recorder demands a value, the owner supplies an admissible value, the device attaches that value to the file, and the file is later copied to a computer that preserves the metadata exactly as given. Years afterward, a listing of that folder can show a modification date decades in the future, not because any component behaved incorrectly, but because a required field was never made to mean anything by the person who filled it in.

The interesting move is not to dismiss the resulting timestamp as garbage, nor to silently correct it once its impossibility is noticed. Both responses discard information. The timestamp is genuine evidence, just not evidence of the event it appears to describe. Properly read, a recording metadata field reading a year that has not yet occurred is excellent evidence about the state of a particular clock inside a particular device at the moment of writing, and poor to worthless evidence about when the recorded sound actually occurred. Four different quantities are being conflated whenever a person treats a timestamp as if it settled the question of when something happened: the event itself, the device's declared account of the event, the time at which some independent party can establish that the resulting record already existed, and the causal order among records that can be established without reference to any clock at all. Keeping these four separate is the discipline this essay is about.

\section{Git: Immutable Objects, Revisable Ancestry}

Version control systems make the same distinction visible in a more consequential setting, because they are widely treated as authoritative historical records and because the metadata in question can be rewritten deliberately and at scale. A Git commit's author and committer dates are ordinary fields inside the commit object, no different in kind from a file's modification time. They can be set to any value at creation, and an existing commit can be amended to carry a new value, which produces a new object with a new hash. Rewriting a commit earlier in a chain forces every descendant to be reconstructed as well, since each commit's identity depends on its parent's hash. A short sequence of ordinary commands, setting an environment variable for the desired date, amending a commit without changing its content, and repeating the operation across a rebase, is sufficient to produce a history in which every commit carries a chosen date rather than the date on which the work actually occurred. A force push then moves the branch reference to point at the reconstructed chain.

What this technique exposes is not a flaw in Git so much as a frequently overlooked structural fact about it. Git's content-addressed objects are genuinely immutable: a commit cannot be edited in place, only superseded by a different object with a different hash. The commit DAG itself does not change when a rewrite occurs; every object that ever existed remains exactly as it was, reachable from wherever it was reachable before. What changes is which history a given branch name currently selects out of that ever-growing graph. A repository's main branch is a pointer, and pointers can be moved. Before a rewrite, an original chain of commits terminates at a reference; a backup branch can preserve a pointer to that same chain; the rewrite constructs a new chain whose rewritten commits acquire new identities from the first changed commit onward, since each commit's hash depends on its content and its parent's hash, and moves the main reference onto this new chain. Nothing about this operation destroys the original commits. They remain reachable through the backup branch, or through any other reference or clone that still points at them. What has changed is which continuation a particular name now recognizes.

This licenses a compact formulation. A commit hash establishes the integrity of the commit's content relative to itself; it does not establish that the commit's declared metadata is historically accurate. A branch establishes a chosen continuation; it does not establish that no alternative continuation ever existed. A force push changes the continuation a remote branch recognizes; it does not travel backward and alter an existing commit, because doing so is impossible by construction. Git therefore holds two properties in tension that sound contradictory until they are separated: objects are immutable, and the history a reference selects is revisable. The resolution is that a history, in the sense that matters for a branch, is not a set of objects but a chosen path through the graph of all objects that currently exist, and that path can be redirected without touching any individual object on it.

A further subtlety compounds the first. Backdating a commit does not establish that the commit existed on the declared date; it only establishes, via the hash, that the commit's content has not changed since it was created, whenever that was. A commit created today with a metadata field claiming three years ago carries a hash that protects the claim from subsequent tampering without certifying that the claim was true when made. Establishing the stronger claim, that a particular object demonstrably existed no later than some date, requires evidence external to the object's own self-declared timestamp: an independently retained copy on a remote server whose access logs are trusted, a signed and dated publication, a transparency log, or some other outside witness. Git alone stores what might be called declared history, the account a commit gives of itself, and structural history, the graph of what preceded what according to the parent pointers currently in force. It provides comparatively little of what might be called witnessed history, independent evidence that some outside observer actually encountered a given object at a given time, unless that evidence is deliberately constructed and kept separate from the repository's own self-description.

The whole of this section is, in the end, one instance of a single short principle, worth stating on its own rather than leaving implicit in the surrounding discussion of hashes and branches.
\[
\boxed{\text{Integrity preserves what was asserted; it does not establish that the assertion was true.}}
\]
A hash, a signature, or a content-addressed identifier can guarantee that a claim has not been altered since it was made. None of these mechanisms can guarantee that the claim was correct when it was made. Confusing the two is what allows a formally tamper-evident record to be mistaken for a formally verified one.

\section{An External Witness Layer}

The natural response to this gap is to add an append-only record that sits outside any single repository and outside any single party's control, so that a proposed rewrite cannot simply replace what came before it in every copy simultaneously. Such a layer need not be a blockchain in the popular sense; a signed and timestamped publication, a Merkle transparency log, or a sufficiently distributed set of independent mirrors can supply closely related guarantees with considerably less machinery. What matters is not the specific mechanism but the property it is meant to supply: that a claim about a commit's existence, once published, becomes itself a historical event with its own evidence trail, independent of whatever any single branch pointer currently asserts.

The correct design for such a layer is easy to get wrong in a way that reproduces the original problem at a larger scale. If the external log is used to decide which Git history is canonical, then the branch pointer's authority has simply been relocated to a different, ostensibly more trustworthy pointer, and the same distinction between a chosen continuation and an established past reappears one level up. The more useful design keeps the layer's role strictly evidentiary: it records that a chain terminating at a given commit was observed at a given time, and later records that a different chain, sharing an earlier ancestor but diverging afterward, was subsequently observed. Both observations persist. Two parties can disagree about which chain ought to be developed further while agreeing, on the strength of the log, that the first chain existed, that a second chain was later proposed, and that some reference eventually moved from the first to the second. Git supplies the identity of objects; the log supplies the timing and existence of observations of those objects; a separate and explicit policy decision supplies which observed history is chosen for continuation. Collapsing these three operations into one is precisely the move that made the original timestamp problem hard to see in the first place.

A second non-implication belongs alongside the first, so that the witness layer is not mistaken for a stronger guarantee than it actually supplies.
\[
\boxed{\text{witnessed existence} \;\not\Rightarrow\; \text{semantic truth}.}
\]
A transparency log that proves a document existed, unaltered, at a given time proves exactly that and nothing further. It does not certify any assertion contained inside the document. The witness layer answers the question of when and whether a claim was made; it has no bearing on whether the claim was correct. Treating a well-witnessed record as thereby a verified one repeats, at the level of the log rather than the repository, the same error the Sony recorder illustrates at the level of a single file.

\section{Predictions as Commitments}

The same structure appears again, in a form with more immediate practical stakes, wherever a system produces an output about a state it has not yet observed. Consider a model that generates an image of the unseen side of a physical object from a single photograph of its visible side, or one that predicts a masked region of a scene from the surrounding context, or one that classifies a photograph. Each of these outputs is a commitment made under a specific, limited evidential horizon: a definite claim, produced by a definite model at a definite time, given exactly the information that was available and no more. When the object is subsequently rotated, when the mask is removed, or when independent evidence establishes that two separately classified photographs in fact depict the same thing, a comparison becomes possible between the original commitment and the later observation.

The temptation at this point is to discard the earlier, now-falsified output and retain only the corrected result, on the reasoning that the mistaken prediction served no further purpose once the truth was known. This discards exactly the information that made the correction possible to evaluate in the first place. The mistaken prediction is not incidental to the experiment; it is the experiment. Without a retained record of what a model committed to before the additional evidence arrived, there is no way to later ask what that model, at that checkpoint, given that evidence and no more, actually inferred. A model that turns out to have been right for the wrong reasons, or confidently wrong, or appropriately uncertain, cannot be distinguished from one another once the original commitment has been overwritten by the corrected answer.

This reframes an otherwise unremarkable observation about model errors, that a system asked about typesetting reliably renders a particular proper name incorrectly regardless of surrounding context, as more than a benchmark failure to be patched in the next release. Such an error, if reproducible, is a standing counterexample: there exists a context in which the model's transformation of its input violates a known and checkable constraint. The counterexample does not need to be corrected out of the historical record once a later version happens to get it right; it needs to be retained as a fixed point against which every later version can continue to be tested, so that an apparent improvement can be distinguished from a coincidental one, and so that a regression can be caught rather than silently reintroduced.

The multimodal case makes the same point at larger scale. A single photograph supports one reconstruction of an occluded region; additional photographs from other angles, a LiDAR scan, and an infrared pass each supply further constraints. The interesting question is not simply whether the final, most-corroborated reconstruction is more accurate than the first, which is close to true by construction, but whether accuracy improves monotonically as genuinely independent evidence accumulates. A system whose reconstructions become worse after receiving an additional, non-redundant observation has a demonstrable defect in how it integrates evidence, and that defect is only detectable if the sequence of earlier reconstructions has been preserved rather than replaced in place by whatever the system currently believes.

One qualification is necessary here. Additional modalities are not automatically additional truth. Several instruments mounted on a common platform can share a single upstream error, such as a mis-estimated pose, and will then agree with one another for a reason that has nothing to do with the world they are jointly measuring. What matters is not the count of modalities but the independence of their provenance: whether the paths by which two observations came to exist share a common point of failure. A useful record of evidence therefore needs to track not just what was observed but how each observation was produced, so that apparent corroboration between two sources can itself be checked for hidden dependence.

\section{Admissible Interpretation Spaces}

These cases share enough structure to admit a common formalization, provided one preliminary distinction is made explicit before any constraint sets are written down. An observation is not itself a constraint on $\Omega$; it becomes one only once an interpretation has been applied to it, and the interpretation is a separate, fallible commitment that deserves its own place in the record rather than being silently absorbed into the observation. Write $O_i$ for the raw observation, $I_i$ for the interpretation map that carries $O_i$ into a claim about the quantity of interest, and $\tau_i$ for the measurement tolerance attached to that claim. The constraint actually entering the formalism is then
\[
C_i \;=\; \Phi(O_i, I_i, \tau_i) \;\subseteq\; \Omega,
\]
and provenance, in the sense this essay has used the word throughout, means retaining $O_i$, $I_i$, and $\tau_i$ separately rather than only their combination $C_i$. This separation matters because it determines where a later contradiction is diagnosed. If only $C_i$ is retained, a contradiction can only be read as ``this evidence is bad.'' If $O_i$ and $I_i$ are retained separately, the same contradiction can instead implicate the interpretation map while leaving the underlying observation untouched, which is exactly the diagnosis the Sony recorder case below requires. The same separation generalizes immediately to machine learning pipelines, where raw sensor readings are rarely the objects a model reasons over directly; preprocessing, registration, coordinate transforms, and semantic classification are all instances of $I_i$, and an empty intersection downstream can originate in any one of them rather than in the sensor itself.

Let $\Omega$ be the space of candidate states of the world relevant to some question, equipped with a reference measure $\mu$. Given evidence accumulated up to time $t$, $E_t = \{O_1,\dots,O_t\}$ together with their interpretations, define the admissible interpretation space
\[
\mathbb{A}(E_t) \;=\; \bigcap_{i \le t} C_i .
\]
When evidence is genuinely appended rather than revised, $E_{t+1} = E_t \cup \{O_{t+1}\}$, intersection alone guarantees
\[
\mathbb{A}(E_{t+1}) \;=\; \mathbb{A}(E_t) \cap C_{t+1} \;\subseteq\; \mathbb{A}(E_t).
\]
This is monotonicity, and it follows from set intersection without any appeal to a learning algorithm; it is a property an evidence-integration architecture can be built to guarantee independently of whatever model happens to sit inside it.

Monotonicity by itself is not enough, because shrinking the admissible space measures the strength of accumulated constraints, not their correctness. A single badly calibrated observation can shrink $\mathbb{A}(E_t)$ dramatically, $\mu(\mathbb{A}(E_{t+1})) \ll \mu(\mathbb{A}(E_t))$, while simultaneously excluding the true state, $\omega^\ast \notin \mathbb{A}(E_{t+1})$. A system should therefore be required to satisfy two separate properties rather than one: the nesting condition above, and consistency, $\omega^\ast \in \mathbb{A}(E_t)$ for all $t$. Consistency stated this way is a soundness criterion, not an operational test, since $\omega^\ast$ is ordinarily exactly what is not known at time $t$; it cannot be checked contemporaneously, only presupposed or later found violated. This is not a defect in the formalization so much as the reason the earlier preservation requirement matters. Because $P_t$ and $\mathbb{A}(E_t)$ are retained rather than overwritten, a later expansion of the evidential horizon can retroactively test a soundness claim that was unverifiable when it was made, by checking whether $\omega^\ast$, now better constrained, would have fallen inside $\mathbb{A}(E_t)$ or not. Soundness in this framework is therefore latent at the time a commitment is made and only becomes operationally testable in arrears, as further evidence narrows the space enough to locate $\omega^\ast$ with confidence, which is itself an argument for preserving early commitments rather than an incidental consequence of doing so.

Distinguishing latent soundness from mere contraction yields a three-way failure taxonomy. Expansion is a violation of the nesting condition, $\mathbb{A}(E_{t+1}) \not\subseteq \mathbb{A}(E_t)$, and indicates that new evidence was not actually treated as a constraint on the old space. Contradiction is the collapse $\mathbb{A}(E_{t+1}) = \varnothing$, and indicates that some pair of observations, or some hidden assumption in one of the interpretation maps connecting them, cannot be jointly satisfied. False exclusion is the case $\omega^\ast \notin \mathbb{A}(E_{t+1})$ while the space remains nonempty, which is the most dangerous failure because it produces a confident, internally consistent, and wrong result with no visible symptom at all, and which can only be identified after the fact, once later evidence permits the retroactive test just described.

The Sony recorder example becomes a clean instance of contradiction once it is put in these terms, and it is exactly the case that motivates keeping $O_i$ and $I_i$ separate rather than working directly with $C_i$. The device supplies an observation, $O_{\text{Sony}}: T_{\text{device clock}} = 2036$, and an interpretation map, $I_{\text{clock}\to\text{event}}: T_{\text{device clock}} \mapsto T_{\text{event}}$, which is applied by default and without comment by anyone who reads a file's timestamp as the time of the event it recorded. Only the combination, $C_{\text{device}} = \Phi(O_{\text{Sony}}, I_{\text{clock}\to\text{event}}, \tau) = \{\omega : T_{\text{event}} = 2036\}$, contradicts independent evidence about when the device was purchased and used, $C_{\text{obs}} = \{\omega : T_{\text{event}} < 2026\}$, giving an empty intersection. A system that retains only $C_{\text{device}}$ can report no more than that the data are bad. A system that retains $O_{\text{Sony}}$ and $I_{\text{clock}\to\text{event}}$ separately can instead ask which of the two is at fault, and in this case the answer is clear: the observation, $T_{\text{device clock}} = 2036$, is perfectly accurate, and the interpretation, that device clock time equals event time, is the unjustified assumption doing the damage. The contradiction, properly investigated, identifies not bad data but a hidden and unwarranted interpretation map, and this diagnosis is only available if the observation and the interpretation were both retained long enough to be examined separately rather than the file's timestamp being corrected on sight.

A companion caution concerns what admissible sets leave out. Two models can make the same categorical error while differing enormously in how confidently they made it, and a purely set-based representation collapses that difference. A model assigning probability $0.51$ to a particular outcome and one assigning probability $0.999999$ to the same, incorrect, outcome have both committed the identical error under the admissibility framework, yet the second has made a far more serious miscalibration. The stronger representation is therefore a pair, $(\mathbb{A}_t, p_t)$, an admissible set recording what remains logically compatible with the evidence, held alongside a probability distribution recording how the system's confidence was actually distributed within that space at the time. Subsequent evidence then evaluates both: the set contracts according to the constraint above, and the distribution is scored against the same proper scoring rules already used for calibration, so that the framework supplements rather than replaces ordinary loss functions.

\section{Preserve, Constrain, Continue}

The formalism above supports a compact statement of the operational principle this essay has been building toward. Three operations need to remain distinct wherever a system produces a claim about a state it cannot fully observe. Preservation retains the original commitment, $P_t$, unmodified once evidence at time $t+1$ arrives. Constraint updates the admissible interpretation space in light of new evidence, $\mathbb{A}_{t+1} = \mathbb{A}_t \cap C(E_{t+1})$, without touching $P_t$ itself. Continuation is the separate decision of which admissible interpretation, or which branch, or which reconstructed history, the system will actually build upon going forward.

\[
\text{Preserve}(P_t), \qquad \text{Constrain}(\mathbb{A}_t, E_{t+1}), \qquad \text{Continue}(H_t).
\]

Collapsing these into a single overwrite operation, replace the old claim with the new one, is efficient and is also precisely what destroys the evidence that made the correction possible to evaluate. The alternative, which this essay proposes as a general principle rather than a domain-specific fix, is: do not correct the historical commitment; preserve the commitment, append the correction, record the discrepancy, and correct the continuation rather than the record.

This formulation resolves the apparent tension between two properties that a good historical system needs simultaneously. History, in the sense of the accumulated record of past commitments and observations, should be monotonic: nothing already recorded is deleted or silently altered by later information. Belief, in the sense of which interpretation is currently treated as most likely or which branch is currently being developed, should remain fully revisable in light of new evidence. And every commitment recorded along the way should remain falsifiable rather than being quietly reinterpreted after the fact to have meant something more defensible than what was actually claimed at the time. A system that satisfies all three, an immutable ledger, a revisable interpretation drawn from it, and commitments left intact for future falsification, can be audited in a way that a system which simply keeps its best current answer cannot.

\section{Media Quines and Provenance Graphs}

The constructive counterpart of this principle is an artifact designed from the outset to carry its own provenance across every transformation applied to it. Call such an artifact a media quine: a work that exists as a text, an audiobook, an abridged explanation, a film, a simulation, or a dataset, where each of these representations is produced from, and retains enough information to be related back to, the others. The relevant analogy is not a lossless encoding, which would require every representation to be reconstructible from any other, but a ring species, in which adjacent forms remain interoperable while forms separated by many intermediate steps may no longer be directly interchangeable, and yet the whole chain remains connected through its sequence of local transformations.

The more accurate structure is a graph rather than a ring. Let $Q = (V, T, P)$, where $V$ is a set of representations, $T$ a set of transformations between them, and $P$ the provenance record attached to each transformation, recording what process produced it, from what inputs, and under what evidential horizon. A manuscript may transform directly into a film as well as into an audiobook that separately transforms into an abridged explainer; a simulation may transform into a dataset that has no direct counterpart among the earlier representations at all. Multiple paths through this graph can connect the same two nodes.

An unqualified requirement that any two paths arriving at the same node agree outright is too strong to be useful, since an audiobook and a simulation of the same underlying work are not expected to be identical, and are not failing at anything when they differ in the respects neither of them promised to preserve. Each transformation $e \in T$ should instead carry, alongside its provenance record, an explicit preservation contract $J_e$, a specification of which invariants of its source the transformation claims to carry forward and which it does not. An audiobook's contract might promise to preserve narrative sequence and characters' stated claims while making no promise about visual detail; a simulation's contract might promise to preserve certain measurable dynamics while making no promise about surface narration at all. Comparison between two paths, $A \xrightarrow{f} B \xrightarrow{g} D$ and $A \xrightarrow{h} C \xrightarrow{k} D'$, is then a test not of whether $D$ and $D'$ coincide but of whether some invariant $J$ that both paths jointly claim to preserve agrees between them,
\[
J(D) \;\stackrel{?}{=}\; J(D'),
\]
and a discrepancy here, $\Delta_{\text{path}} = d\bigl(J(D), J(D')\bigr)$, is informative in exactly the way a discrepancy between a prediction and a later observation is informative elsewhere in this essay. It identifies a place where two independently derived representations fail to agree on a quantity both were supposed to hold fixed, which implicates at least one transformation along one of the two paths, whether or not either output is otherwise defective. A media quine, on this formulation, is not a claim that all of its representations agree; it is a claim that every transformation states explicitly what it preserves, and that those particular, stated claims remain falsifiable against every other path that makes the same claim.

This gives the media quine a stronger property than mere translation between formats. Because each transformation retains its own provenance, later scrutiny does not need to trust a single canonical representation; it can compare independently derived representations against one another and treat their agreement or disagreement as evidence in its own right, in the same way that independent corroboration across modalities strengthens or weakens a reconstruction in the earlier discussion, and for the same reason: what matters is not how many representations exist but how independently they were produced.

\section{A Note on Declared Provenance in Practice}

The distinction developed here has a mundane but recurring practical consequence outside of research contexts. A document, a repository, or a filing that asserts ownership, authorship, or evidentiary status inside its own fields, an owner field naming a party, a severity field assigning a word like ``high,'' a chain-of-custody block declaring a hash and a controller, has recorded a claim with the same epistemic status as a device's self-reported clock. Such fields can be produced with complete internal consistency, cross-referenced, versioned, and hashed, and none of that machinery establishes that the claims are true; it establishes only that the claims, whatever they are, have not been altered since they were recorded. Systems that present formal-sounding provenance apparatus, version numbers, chain-of-custody sections, cryptographic hash placeholders, around declarations that are never checked against anything external are exhibiting the identical structure as the Sony recorder: an evidentiary vocabulary attached to a value nobody has actually witnessed. The remedy is the same in both cases. The declaration is retained as data about what was claimed and by whom, and it is kept separate from any inference about what was actually the case, which requires independent witnessing that the declaration alone cannot supply.

\section{Conclusion}

The claim this essay has argued for is narrower than it might first appear, and that narrowness is what makes it checkable across such different domains. A record's self-description is not a substitute for an outside witness to that record, whether the record is a timestamp, a commit, a model's prediction, or a legal filing's own provenance section. Later evidence should be allowed to narrow the space of interpretations that remain admissible, and to attach a discrepancy to whatever earlier commitment it falsifies, without ever being permitted to reach backward and rewrite what that commitment actually was. The Sony recorder's impossible date, the Git commit with a manufactured timestamp, and the vision model's confident guess about an object's unseen side are the same structure viewed at three different scales, and the same discipline, preserve the commitment, constrain the interpretation, correct the continuation, applies to each without modification.

What this essay has actually been describing is not, in the end, immutability as such. Append-only logs, content-addressed stores, and tamper-evident records are established engineering, and none of them is the paper's contribution. The contribution is a separation of three properties that are routinely collapsed into one system, one field, or one overwrite operation, and that need not be collapsed at all.
\[
\boxed{
\begin{array}{c}
\text{preservation without correction} \\
+\;\text{constraint without erasure} \\
+\;\text{continuation without canonicalization}
\end{array}}
\]
A claim can remain permanently intact in the record; be demonstrably false in light of everything now known; and cease to govern what the system builds on next, all three at once, without contradiction, provided the architecture keeps these as three separate operations rather than one. The Sony file, the rewritten commit, and the falsified prediction are none of them, in this framework, mistakes to be cleaned out of the historical record. They are the record functioning exactly as it should.

\end{document}

